\appendix
\chapter{Application Details}
\label{app:application}

\section{Data Organization}

Fracture examples are stored as static OBJ files organized in a hierarchical directory structure:

\begin{verbatim}
meshes/
  two_pieces/
    <category>/
      <mesh_hash>/
        fractured_0/
          piece_0.obj
          piece_1.obj
  multi_pieces/
    <category>/
      <mesh_hash>/
        fractured_0/
          piece_0.obj
          piece_1.obj
          piece_2.obj
\end{verbatim}

At startup, the Django application automatically discovers all valid fractures by walking this directory tree. Each fracture is assigned a deterministic 12-character ID derived from its path, ensuring stable URLs across restarts. The discovery logic filters for directories starting with \texttt{fractured\_} and containing at least two \texttt{piece\_*.obj} files.

\section{API Endpoints}

The Django server exposes three endpoints:

\begin{itemize}
    \item \texttt{GET /} --- Renders the main page with fracture lists injected into the template context.
    \item \texttt{GET /api/fracture/<id>/} --- Returns JSON metadata and static OBJ URLs for a specific fracture.
    \item \texttt{POST /api/predict/} --- Accepts a fracture ID and model variant, proxies to the GPU server or returns cached results.
\end{itemize}

The FastAPI inference server exposes:

\begin{itemize}
    \item \texttt{GET /health} --- Returns server status and CUDA availability.
    \item \texttt{GET /models} --- Lists all six available model variants with descriptions and maximum piece counts.
    \item \texttt{POST /predict} --- Accepts OBJ files as multipart uploads, runs the full Jigsaw inference pipeline (preprocessing, forward pass, global alignment), and returns the assembled meshes as OBJ text.
\end{itemize}

The server validates all inputs: it rejects requests with fewer than 2 pieces, more pieces than the selected model supports (2 for two-piece models, 4 for multi-piece), and files that are not valid UTF-8 text.

\section{User Interface Details}

The application presents a single-page layout consisting of two main regions: a scrollable narrative column on the left and a sticky interactive demo panel on the right.

\subsection{Narrative Slides}

The left column contains five full-viewport-height slides summarizing the thesis:

\begin{enumerate}
    \item \textbf{Problem introduction:} Fracture assembly as a 6-DoF pose estimation problem, applications in archaeology and forensics.
    \item \textbf{Jigsaw pipeline:} Overview of the four-stage architecture (PointNet++, attention, matching, alignment).
    \item \textbf{Proposed improvements:} Pair geometric attention and gated double attention mechanisms.
    \item \textbf{Experimental results:} Key findings on dataset size, architectural modifications, and ICP refinement.
    \item \textbf{Future work:} Scaling to more pieces, loss scheduling, sim-to-real transfer.
\end{enumerate}

Users navigate between slides using arrow buttons or by scrolling.

\subsection{Interactive Demo Panel}

The demo panel remains fixed on the right side of the screen and contains:

\begin{itemize}
    \item \textbf{Model selector:} A dropdown menu listing all six model variants covering different architectural configurations and piece counts.
    \item \textbf{Category filter:} Filters the available examples by object category (e.g., BeerBottle, Cup).
    \item \textbf{Example navigator:} Previous/next buttons and a counter to browse through examples within the selected category.
    \item \textbf{Input viewer:} A 3D viewer displaying the broken fragments in a scattered configuration.
    \item \textbf{Fix Object button:} Triggers inference and displays a loading indicator during processing.
    \item \textbf{Output viewer:} A 3D viewer displaying the assembled result after inference completes.
\end{itemize}

An \textbf{Expand Demo} button opens a fullscreen overlay with a larger 3D viewer and two additional tabs: \textbf{Browse Examples} (a grid of all available fractures with thumbnail cards) and \textbf{Compare Performance} (a side-by-side comparison interface).

\section{User Manual}

\subsection{Step 1: Launch and Overview}

Upon opening the application, the user is presented with a hero section introducing the thesis topic. Scrolling down or clicking \textbf{Explore the Demo} reveals the main content area with narrative slides on the left and the interactive demo on the right.

\begin{figure}[H]
    \centering
    \includegraphics[width=0.75\linewidth]{figures/hero_landing_page.png}
    \caption{Landing page with hero section and call-to-action button.}
    \label{fig:app_hero}
\end{figure}

\subsection{Step 2: Select a Model and View Fragments}

The user selects a model variant from the dropdown (e.g., ``Baseline (2 pieces)'' or ``Double attention (2-4 pieces)''). The category filter narrows the list of available examples, and the previous/next buttons or the fullscreen grid allow browsing. After selecting an example, the input viewer loads and displays the fragment meshes in a scattered arrangement with distinct colors, allowing inspection via rotation, zoom, and pan.

\subsection{Step 3: Run Inference}

Clicking the \textbf{Fix Object} button triggers the inference process. The button displays a loading animation while the request is processed. The system either performs live inference on the GPU server or retrieves a cached result. A status message indicates the source of the result (live or cached).

\begin{figure}[H]
    \centering
    \includegraphics[width=0.75\linewidth]{figures/app_result.png}
    \caption{Assembled result displayed in the output viewer after clicking Fix Object. The status bar indicates whether the result was computed live or retrieved from cache.}
    \label{fig:app_result}
\end{figure}

\subsection{Step 4: Inspect and Compare}

The output viewer displays the assembled object. The user can rotate and zoom to inspect alignment quality (Figure~\ref{fig:inspect_result}). Switching to a different model variant and clicking Fix Object again allows comparing reconstruction quality across architectures. The \textbf{Expand Demo} button opens a fullscreen overlay with a larger viewer and a \textbf{Browse Examples} grid for quick selection (Figure~\ref{fig:app_fullscreen}).

\begin{figure}[H]
    \centering
    \begin{subfigure}[t]{0.48\linewidth}
        \centering
        \includegraphics[width=\linewidth]{figures/show_zoom.png}
        \caption{Zoomed view of the assembled object.}
    \end{subfigure}
    \hfill
    \begin{subfigure}[t]{0.48\linewidth}
        \centering
        \includegraphics[width=\linewidth]{figures/show_zoom_rotation.png}
        \caption{Close-up inspection of the fracture pattern.}
    \end{subfigure}
    \caption{Detailed inspection of the reconstruction result from different viewing angles.}
    \label{fig:inspect_result}
\end{figure}

\begin{figure}[H]
    \centering
    \includegraphics[width=0.75\linewidth]{figures/app_fullscreen.png}
    \caption{Fullscreen demo view with the Browse Examples grid for quick selection.}
    \label{fig:app_fullscreen}
\end{figure}

\section{Functional Requirements}

The application satisfies the following functional requirements:

\begin{itemize}
    \item \textbf{F1 -- Fracture visualization:} Display broken 3D fragments in an interactive viewer with rotation, zoom, and pan controls.
    \item \textbf{F2 -- Model selection:} Allow the user to choose between six model variants covering different architectural configurations and piece counts.
    \item \textbf{F3 -- Inference execution:} Perform end-to-end fracture assembly inference (feature extraction, segmentation, matching, and global alignment) given a set of OBJ fragments.
    \item \textbf{F4 -- Result visualization:} Display the assembled object in a separate 3D viewer for inspection.
    \item \textbf{F5 -- Graceful degradation:} Operate without a GPU server by serving precomputed results from a static cache.
    \item \textbf{F6 -- Example browsing:} Provide navigation controls and a fullscreen grid for browsing the complete set of available examples.
\end{itemize}

\section{Non-Functional Requirements}

\begin{itemize}
    \item \textbf{Performance:} Live inference completes within 120 seconds (the proxy timeout). Precomputed fallback responses are served in under 100ms.
    \item \textbf{Portability:} The Django server runs on any Linux environment with Python 3.12+. The inference server requires CUDA-capable hardware with at least 8 GB VRAM.
    \item \textbf{Maintainability:} The codebase is modular: data processing, model components, training infrastructure, web application, and inference server are clearly separated. Configuration is externalized through environment variables and YAML files.
    \item \textbf{Security:} CSRF protection is enabled on all POST endpoints. Session cookies are configured with \texttt{HttpOnly} and \texttt{SameSite=Lax} flags. The inference server restricts CORS origins via environment configuration.
\end{itemize}
